%!TEX root = ../TM3.tex

\documentclass[
	% pdftex,
	paper=a4,
    numbers=noendperiod,
	% footsepline,
	% headsepline,1
	% oneside,
	twoside,
	open=right, 
	% twocolumn=true,
	BCOR=5mm,
	% DIV=16,
	headinclude=true,
	footinclude=false,
	% headings=optiontohead,
	headings=optiontotoc,
	% headings=optiontoheadandtoc,
	% headlines = 1.25,
	% headheight=10mm,
	% footlines = 1.25,
	% footheight=10mm,
	pagesize,
	% draft=true3
	cleardoublepage=empty,
	% titlepage=true,
	titlepage=firstiscover,
	% abstract=true,
	% toc=bibnumbered,
	% toc=listofnumbered,
	chapterprefix=true,
	appendixprefix=true,
	% 11pt,
	ngerman,
	% scrheadings,
	% listof=totoc,
  	% bibliography=totoc,
	]{scrbook}
% \areaset[5mm]{160mm}{260mm}
% \setlength{\marginparwidth}{1.0\marginparwidth}
\usepackage[onehalfspacing]{setspace}
% \usepackage[doublespacing]{setspace}
\KOMAoptions{DIV=last}

\usepackage[
	% hidelinks,
	colorlinks=false,
	linkcolor=blueLink,
	% urlcolor=blueLink,
	% citecolor = black,
	]{hyperref}
\hypersetup{
  pdftitle    = {Technische Mechanik 3},
  pdfsubject  = {Grundlagen der Kinematik und Kinetik},
  pdfauthor   = {Lukas Moschen},
  % pdfkeywords = {Stichwort1, Stichwort2 ...} ,
  % pdfcreator  = {Mit welcher Anwendung i.d.R. pdflatex},
  % pdfproducer = {LaTeX with hyperref}
}

%header and footer
\usepackage{scrlayer-scrpage}
 %  \renewcommand*{\chaptermarkformat}{%
 %    \thechapter\autodot\enskip%
 %  }
	% \automark[chapter]{section}
	% \pagestyle{scrheadings}
	% % \clearscrheadfoot
	% \ihead[]{\headmark}
	% \chead[]{}
	% \ohead[]{}
	% \ifoot[]{}
	% \cfoot[\pagemark]{\pagemark}
	% \ofoot[]{}

\newcommand{\MONTH}{%
  \ifcase\the\month
  \or Januar% 1
  \or Februar% 2
  \or März% 3
  \or April% 4
  \or Mai% 5
  \or Juni% 6
  \or Juli% 7
  \or August% 8
  \or September% 9
  \or Oktober% 10
  \or November% 11
  \or Dezember% 12
  \fi}

\usepackage[
	% subpreambles=true
]{standalone}


\usepackage[
			hyperref=true,
            url=false,
            isbn=false,
            doi=false,
            % backref=true,
            style=authoryear,
            % citereset=section,
            maxcitenames=2,
            maxbibnames=100,
            backend=bibtex,
            % defernumbers=true,
            bibencoding=utf8,
            % block=none,
            autolang = other,
            natbib = true,
            uniquename=init,
            giveninits = true,
            ]{biblatex}
% \DeclareNameAlias{sortname}{first-last}
% \renewcommand*{\bibfont}{\small}
\addbibresource{chapters/bibliography.bib}


\usepackage[utf8]{inputenc}
\usepackage[T1]{fontenc}
\usepackage[ngerman]{babel}
\usepackage{csquotes}
\usepackage{lmodern}
\usepackage{xcolor}
\usepackage{graphicx}
\graphicspath{{D:/SynologyDrive/Work/Teaching/MCI/Mechatronik_BB_TM3/2020WS/06_figures/sketches/}}
\usepackage{tikz}
\usepackage{pgf}
\usepackage[shortlabels]{enumitem} 
\usepackage[notref,final]{showkeys}
\usepackage{placeins}
\usepackage{supertabular}
\usepackage{transparent}
\usepackage{epigraph}
\usepackage{mathtools}
\usepackage{amsmath,amsfonts,amssymb}
% \usepackage{amsthm}
\usepackage{commath}
\usepackage{bm}
\usepackage{bbm}
\usepackage[d]{esvect}
\usepackage{siunitx}
	\sisetup{
	locale = DE,
	per-mode = symbol,
	output-decimal-marker = {,},
	round-mode=places,
	round-precision=2,
	% round-integer-to-decimal,
	}
% \mathtoolsset{showonlyrefs} 
% \usepackage{relsize}
\usepackage{microtype}
% \usepackage[hang]{footmisc}
% \renewcommand{\footnotemargin}{1em}
\interfootnotelinepenalty=10000
% \usepackage{booktabs}
\usepackage{pgfplots}
\pgfplotsset{compat=1.16}
\usepgfplotslibrary{external}
% \usetikzlibrary{external}
%   % \tikzset{external/export=false}
	  % \tikzset{external/up to date check=simple}
	% \tikzset{external/force remake = true}
	\tikzset{external/figure name=Fig_}
	\tikzexternalize[prefix=figures/]
\usepackage{pgfplotstable}
\usepgfplotslibrary{groupplots}
\usepackage{empheq}
\usepackage{pgfpages}
\usepackage{xparse}
\usepackage{units}
\usepackage[figuresright]{rotating}
\usepackage{adjustbox}
\usepackage{pifont}
\usepackage{cancel}
\usepackage{datetime}
\usepackage{ifthen}
\newcommand{\sw}{0}
\usepackage{nth}

% DEFINE COLORS:
\ifthenelse{\equal{\sw}{0}}
{
	\definecolor{myBlue}{RGB}{0,73,131}
	\definecolor{myOrange}{RGB}{244,155,0}
	\definecolor{myGreen}{RGB}{128,128,0}
	\definecolor{myRed}{RGB}{170,0,0}
	\definecolor{blueLink}{RGB}{0,73,131}
	\colorlet{lightBlue}{myBlue!50}
}
{
	\definecolor{myBlue}{RGB}{85,85,85}
	\definecolor{myOrange}{RGB}{198,198,198}
	\definecolor{myGreen}{RGB}{55,55,55}
	\definecolor{myRed}{RGB}{33,33,33}
	\definecolor{blueLink}{RGB}{0,0,0}
	\colorlet{lightBlue}{myBlue!50}
}

% % FORMAT CAPTIONS:
% \usepackage[labelfont={color=blueLink,bf},format=hang]{caption}
% \captionsetup[figure]{name=Abb.}
% \captionsetup[table]{name=Tab.}


% % PARAGRAPH INDENT AND SKIP:
% \setlength{\parindent}{0mm}
% \setlength{\parskip}{3mm}
% \setlength{\fboxsep}{2mm} % margin box in math!

% RE-NEW COMMANDS
\let\oldsin\sin
\let\oldcos\cos
\let\oldtan\tan
\let\oldcot\cot
\let\oldarcsin\arcsin
\let\oldarccos\arccos
\let\oldarctan\arctan
\DeclareMathOperator{\arccot}{arccot}
\let\oldarccot\arccot
\let\oldRe\Re
\let\oldIm\Im
\let\oldmax\max
\let\oldmin\min
\renewcommand{\sin}[1]{\oldsin\left(#1\right)}
\renewcommand{\cos}[1]{\oldcos\left(#1\right)}
\renewcommand{\tan}[1]{\oldtan\left(#1\right)}
\renewcommand{\cot}[1]{\oldcot\left(#1\right)}
\renewcommand{\arcsin}[1]{\oldarcsin\left(#1\right)}
\renewcommand{\arccos}[1]{\oldarccos\left(#1\right)}
\renewcommand{\arctan}[1]{\oldarctan\left(#1\right)}
\renewcommand{\arccot}[1]{\oldarccot\left(#1\right)}
\renewcommand{\Re}[1]{\oldRe\left(#1\right)}
\renewcommand{\Im}[1]{\oldIm\left(#1\right)}
\renewcommand{\max}[1]{\oldmax\left(#1\right)}
\renewcommand{\min}[1]{\oldmin\left(#1\right)}

% NEW COMMANDS:
\newcommand{\prsn}[1]{\textsc{#1}}
\newcommand{\mtrx}[1]{\bm{\MakeUppercase{#1}}}
\newcommand{\vctr}[1]{\vv{{#1}}}
\newcommand{\uvctr}{\vv{{e}}}
\newcommand{\euler}{\mathrm{e}}
\newcommand{\imag}{\mathbbm{i}}
% \newcommand{\vctr}[1]{\vv{\bm{#1}}}
% \newcommand{\uvctr}{\vv{\bm{e}}}
% \newcommand{\bs}[1]{\boldsymbol{#1}}
% \newcommand{\dtq}{\dot{q}}
% \newcommand{\dr}{\dot{\bs{r}}}
% \newcommand{\pa}{\partial}
% \newcommand{\vp}{\varphi}
% \newcommand{\ve}{\varepsilon}
% \newcommand{\dtp}{\dot{\varphi}}
% \newcommand{\ddtp}{\ddot{\varphi}}
% \newcommand{\myfrac}[2]{\frac{\displaystyle #1}{\displaystyle #2}}
% \newcommand{\ol}[1]{\overline{#1}}
% \newcommand*\circled[1]{\tikz[baseline=(char.base)]{
% 		\node[shape=circle,draw,inner sep=1.8pt] (char) {#1};}}
% \newcommand{\pnr}{\underset{+}{\mathlarger{\mathlarger{\rightarrow}}}}
% \newcommand{\pnl}{\underset{+}{\mathlarger{\mathlarger{\leftarrow}}}}
% \newcommand{\pqo}{\underset{+}{\mathlarger{\mathlarger{\uparrow}}}}
% \newcommand{\pqu}{\underset{+}{\mathlarger{\mathlarger{\downarrow}}}}
% \newcommand{\pmr}{\underset{+}{\mathlarger{\mathlarger{\rotatebox[origin=c]{90}{$\curvearrowleft$}}}}}
% \newcommand{\pml}{\underset{+}{\mathlarger{\mathlarger{\rotatebox[origin=c]{90}{$\curvearrowright$}}}}}
% \newcommand{\cmark}{\ding{51}}%
% \newcommand{\xmark}{\ding{55}}%
% \newcommand{\myBox}[1]{\fcolorbox{myBlue}{white}{#1}}
% % underline math mode:
% \newsavebox\MBox
% \newcommand\Cline[2][red]{{\sbox\MBox{$#2$}%
% 		\rlap{\usebox\MBox}\color{#1}\rule[-1.2\dp\MBox]{\wd\MBox}{0.5pt}}}

% \usepackage{framed}
% \renewenvironment{leftbar}[1][\hsize]
% {% 
% 	\def\FrameCommand 
% 	{%	
% 		{\color{lightBlue}\vrule width 2pt}%
% 		\hspace{4pt}
% 		%\fboxsep=\FrameSep\colorbox{lightgray}%
% 	}%
% 	\MakeFramed{\hsize#1\advance\hsize-\width\FrameRestore}%
% }
% {\endMakeFramed}
% \newcounter{example}[chapter]
% \newcounter{prevequation}
% \NewDocumentEnvironment {example} {}
% {%
% 	\refstepcounter{example}%
% 	\setcounter{prevequation}{\value{equation}}%
% 	\setcounter{equation}{0}%
% 	\def\theequation{\alph{equation}}%
% 	\def\theHequation{\thechapter.\theexample.\alph{equation}}%
% 	\leftbar
% 	\textcolor{blueLink}{\textbf{Beispiel \thechapter.\theexample}}\\[1mm]\nopagebreak\ignorespaces
% }
% {%
% 	\endleftbar
% 	\setcounter{equation}{\value{prevequation}}%
% 	\ignorespacesafterend
% }

\usepackage[noabbrev]{cleveref}

\let\oldFootnote\footnote
\newcommand\nextToken\relax

\renewcommand\footnote[1]{%
    \oldFootnote{#1}\futurelet\nextToken\isFootnote}

\newcommand\isFootnote{%
    \ifx\footnote\nextToken\textsuperscript{,}\fi}